% Options for packages loaded elsewhere
\PassOptionsToPackage{unicode}{hyperref}
\PassOptionsToPackage{hyphens}{url}
\documentclass[
]{article}
\usepackage{xcolor}
\usepackage{amsmath,amssymb}
\setcounter{secnumdepth}{-\maxdimen} % remove section numbering
\usepackage{iftex}
\ifPDFTeX
  \usepackage[T1]{fontenc}
  \usepackage[utf8]{inputenc}
  \usepackage{textcomp} % provide euro and other symbols
\else % if luatex or xetex
  \usepackage{unicode-math} % this also loads fontspec
  \defaultfontfeatures{Scale=MatchLowercase}
  \defaultfontfeatures[\rmfamily]{Ligatures=TeX,Scale=1}
\fi
\usepackage{lmodern}
\ifPDFTeX\else
  % xetex/luatex font selection
\fi
% Use upquote if available, for straight quotes in verbatim environments
\IfFileExists{upquote.sty}{\usepackage{upquote}}{}
\IfFileExists{microtype.sty}{% use microtype if available
  \usepackage[]{microtype}
  \UseMicrotypeSet[protrusion]{basicmath} % disable protrusion for tt fonts
}{}
\makeatletter
\@ifundefined{KOMAClassName}{% if non-KOMA class
  \IfFileExists{parskip.sty}{%
    \usepackage{parskip}
  }{% else
    \setlength{\parindent}{0pt}
    \setlength{\parskip}{6pt plus 2pt minus 1pt}}
}{% if KOMA class
  \KOMAoptions{parskip=half}}
\makeatother
\setlength{\emergencystretch}{3em} % prevent overfull lines
\providecommand{\tightlist}{%
  \setlength{\itemsep}{0pt}\setlength{\parskip}{0pt}}
\usepackage{bookmark}
\IfFileExists{xurl.sty}{\usepackage{xurl}}{} % add URL line breaks if available
\urlstyle{same}
% fallback for those not using the hyperref driver hyperxmp:
\makeatletter
\@ifundefined{xmpquote}{\newcommand{\xmpquote}[1]{#1}}{}
\makeatother
\hypersetup{
  hidelinks,
  pdfcreator={LaTeX via pandoc}}

\author{}
\date{}

\begin{document}

\section{Testing with input space
partitioning}\label{testing-with-input-space-partitioning}

We are testing a \texttt{csv2json} program written in Rust that we found on 
\href{https://github.com/benkay86/csv2json}{GitHub}
\subsection{Specification of the Program Under Test}
\texttt{CSV} (Comma Separated
Values) files are used for representing data in spreadsheet programs.
According to
\href{https://www.rfc-editor.org/info/rfc4180/\#section-2}{RFC 4180}: 1.
Each record is located on a separate line, delimited by a line break
(CRLF). For example:

\begin{quote}
\begin{verbatim}
 aaa,bbb,ccc CRLF
 zzz,yyy,xxx CRLF
\end{verbatim}
\end{quote}

\begin{enumerate}
\def\labelenumi{\arabic{enumi}.}
\setcounter{enumi}{1}
\tightlist
\item
  The last record in the file may or may not have an ending line break.
  For example:
\end{enumerate}

\begin{quote}
\begin{verbatim}
  aaa,bbb,ccc CRLF
  zzz,yyy,xxx
\end{verbatim}
\end{quote}

\begin{enumerate}
\def\labelenumi{\arabic{enumi}.}
\setcounter{enumi}{2}
\tightlist
\item
  There maybe an optional header line appearing as the first line of the
  file with the same format as normal record lines. This header will
  contain names corresponding to the fields in the file and should
  contain the same number of fields as the records in the rest of the
  file (the presence or absence of the header line should be indicated
  via the optional ``header'' parameter of this MIME type). For example:
\end{enumerate}

\begin{quote}
\begin{verbatim}
  field_name,field_name,field_name CRLF
  aaa,bbb,ccc CRLF
  zzz,yyy,xxx CRLF
\end{verbatim}
\end{quote}

\begin{enumerate}
\def\labelenumi{\arabic{enumi}.}
\setcounter{enumi}{3}
\tightlist
\item
  Within the header and each record, there may be one or more fields,
  separated by commas. Each line should contain the same number of
  fields throughout the file. Spaces are considered part of a field and
  should not be ignored. The last field in the record must not be
  followed by a comma. For example:
\end{enumerate}

\begin{quote}
\begin{verbatim}
  aaa,bbb,ccc
\end{verbatim}
\end{quote}

\begin{enumerate}
\def\labelenumi{\arabic{enumi}.}
\setcounter{enumi}{4}
\tightlist
\item
  Each field may or may not be enclosed in double quotes (however some
  programs, such as Microsoft Excel, do not use double quotes at all).
  If fields are not enclosed with double quotes, then double quotes may
  not appear inside the fields. For example:
\end{enumerate}

\begin{quote}
\begin{verbatim}
  "aaa","bbb","ccc" CRLF
  zzz,yyy,xxx
\end{verbatim}
\end{quote}

\begin{enumerate}
\def\labelenumi{\arabic{enumi}.}
\setcounter{enumi}{5}
\tightlist
\item
  Fields containing line breaks (CRLF), double quotes, and commas should
  be enclosed in double-quotes. For example:
\end{enumerate}

\begin{quote}
\begin{verbatim}
  "aaa","b CRLF
  bb","ccc" CRLF
  zzz,yyy,xxx
\end{verbatim}
\end{quote}

\begin{enumerate}
\def\labelenumi{\arabic{enumi}.}
\setcounter{enumi}{6}
\tightlist
\item
  If double-quotes are used to enclose fields, then a double-quote
  appearing inside a field must be escaped by preceding it with another
  double quote. For example:
\end{enumerate}

\begin{quote}
\begin{verbatim}
  "aaa","b""bb","ccc"
\end{verbatim}

\mbox{}%
\subsection{}\label{section}

\texttt{JSON} (JavaScript Object Notation) files are a commonly used
serialized data exchange format over the internet.
\href{https://www.rfc-editor.org/info/rfc8259/\#section-1}{RFC 8259}
defines it as: A JSON text is a sequence of tokens. The set of tokens
includes six structural characters, strings, numbers, and three literal
names.
\end{quote}

These are the six structural characters:

\begin{verbatim}
  begin-array     = ws %x5B ws  ; [ left square bracket

  begin-object    = ws %x7B ws  ; { left curly bracket

  end-array       = ws %x5D ws  ; ] right square bracket

  end-object      = ws %x7D ws  ; } right curly bracket

  name-separator  = ws %x3A ws  ; : colon

  value-separator = ws %x2C ws  ; , comma
\end{verbatim}

A JSON value MUST be an object, array, number, or string, or one of the
following three literal names:

\begin{verbatim}
  false
  null
  true
\end{verbatim}

\subsubsection{Input and Output relationship + Corner
Cases}\label{input-and-output-relationship-corner-cases}

In addition to typical inputs, csv2json could encounter corner cases in
the input that affect the JSON output, such as:

\begin{verbatim}
  Empty file input: An Empty CSV file or files with only headers
  Missing/Extra fields: More or less fields than the header specified
  Empty values: Empty fields, partially or fully empty row
  Array style headers: header such as name.first, name.last
\end{verbatim}

\subsection{Input Space Partitioning}\label{input-space-partitioning}

\subsubsection{Component and input
identification}\label{component-and-input-identification}

\texttt{csv2json} takes in a \texttt{CSV} (Comma Separated Value) file
and optional flags and options:

\begin{verbatim}
$ csv2json --in <csv file> > <json file>
\end{verbatim}

and outputs a formatted \texttt{JSON} file. The parameters are as
follows:

\begin{verbatim}
--delimiter (-d): Changes the CSV splitting character.

--dimensional-separator (-D): Creates nested JSON objects based on header syntax.

--arrays (-a): Breaks comma-separated string values into JSON arrays.

--remove-empty-strings: Strips empty string keys.

--remove-empty-objects: Strips empty JSON nodes.

--fold (-F): Folds row-by-row objects into a single object of arrays.

--jsonl: Outputs newline-delimited JSON Lines instead of a standard JSON array.
\end{verbatim}

\subsubsection{Characteristics for partitioning the input
space}\label{characteristics-for-partitioning-the-input-space}

To model this for combinatorial testing, we first partition the input
domain into characteristics and mutually exclusive blocks (levels).

\paragraph{\texorpdfstring{\emph{Characteristic 1: Input File
State}}{Characteristic 1: Input File State}}\label{characteristic-1-input-file-state}

Blocks:
\begin{itemize}
  \item Valid CSV Data
  \item Empty File
  \item Malformed CSV (e.g., unmatched
  quotes).
  \item Header only, missing file
\end{itemize}
\paragraph{\texorpdfstring{\emph{Characteristic 2: Delimiter Option
(\texttt{-\/-delimiter},
\texttt{-d})}}{Characteristic 2: Delimiter Option (-\/-delimiter,
\itemd)}}\label{characteristic-2-delimiter-option---delimiter--d}

Blocks:
\begin{itemize}
\item Default (Comma)
\item Custom (e.g., \texttt{‘;’} or \texttt{‘:’}).
\end{itemize}

\paragraph{\texorpdfstring{\emph{Characteristic 3: Dimensional Separator
(\texttt{-D})}}{Characteristic 3: Dimensional Separator (-D)}}\label{characteristic-3-dimensional-separator--d}

Blocks:
\begin{itemize}
  \item None (Default)
  \item Custom (e.g., Dot .)
  \item Multidimensional (ex. \texttt{a.b.c})
  \item Incorrect formats (\texttt{.a}, \texttt{a.}, \texttt{a..b})
  \item Conflicting keys (\texttt{a} \& \texttt{a.b} in same header)
\end{itemize}

\paragraph{\texorpdfstring{\emph{Characteristic 4: Array Processing
(\texttt{-\/-arrays} or \texttt{-a} with \texttt{-D} to break items to
arrays)}}{Characteristic 4: Array Processing (-\/-arrays or
\itema with
\itemD to break items to arrays)}}\label{characteristic-4-array-processing---arrays-or--a-with--d-to-break-items-to-arrays}

Blocks: Enabled, Disabled.

\paragraph{\texorpdfstring{\emph{Characteristic 5: Output Formatting
(Mutually
Exclusive)}}{Characteristic 5: Output Formatting (Mutually Exclusive)}}\label{characteristic-5-output-formatting-mutually-exclusive}

Blocks:
\begin{itemize}
  \item Standard JSON
  \item Folded (\texttt{-F})
  \item JSONL
  (\texttt{-\/-jsonl})
\end{itemize}

\paragraph{\texorpdfstring{\emph{Characteristic 6: Remove empty
strings/objects (\texttt{-\/-remove-empty-strings},
\texttt{–-remove-empty-objects})}}{Characteristic 6: Remove empty strings/objects (-\/-remove-empty-strings,
\item--remove-empty-objects)}}\label{characteristic-6-remove-empty-stringsobjects---remove-empty-strings--remove-empty-objects}

Blocks:
\begin{itemize}
  \item Use remove empty strings only
  \item Use remove empty blocks only
  \item Use none
  \item Use both Both commands happen to both arrays and objects
  (warned there may be undesired effects) Empty rows, partially empty rows
\end{itemize}

\paragraph{\texorpdfstring{\emph{Characteristic 7: output destination
(\texttt{-\/-out-dir},
\texttt{-–out-name})}}{Characteristic 7: output destination (-\/-out-dir,
\item\/--out-name)}}\label{characteristic-7-output-destination---out-dir--out-name}

Blocks:
\begin{itemize}
  \item Default (stdout)
  \item \texttt{–out-dir},
  \item \texttt{–out-dir} with \texttt{–out-name}
\end{itemize}

\paragraph{\texorpdfstring{\emph{Characteristic 8: Column type
conversion (\texttt{-\/-boolean},
\texttt{–numeric})}}{Characteristic 8: Column type conversion (-\/-boolean,
\item-numeric)}}\label{characteristic-8-column-type-conversion---boolean-numeric}

Blocks:
\begin{itemize}
  \item None
  \item \texttt{-\/-numeric} only
  \item \texttt{-–boolean}
  \item both
\end{itemize}
\end{document}
